\documentclass{leveldoc}
\DocUses{foundations,first_bend,slingshot,around_the_wall,two_stages,thomson_crt}
\DocFingerprint{4702ed9b99bdeaba}

\begin{document}

\section{What the level asks}\label{sec:asks}

The game's text:
\leveldescription

Six shots leave $(\val{xLaunch}, \val{yLaunch})$. Three have kinetic energy
$T = \val{Tslow}$ and three have $T = \val{Tfast}$, in the ratio $\val{ratioT}$.
Each energy is launched at $\val{angSlowDown}^\circ$, $\val{angSlowAxis}^\circ$
and $\val{angSlowUp}^\circ$ to $+x$: the directions stored in the file. The
particle is the same on every shot, $q = \val{q}$ and $m = \val{m}$. There is
no direction cone and no kinetic window.

One gate stands in front of the spots,
$\val{xGlo} \le x \le \val{xGhi}$, $\val{yGlo} \le y \le \val{yGhi}$, with no
condition on the state at the entry. The detector counts only after that entry
(\xeqref{two_stages}{eq:gate}). The three slow shots are counted in the upper
spot $\val{xSlowLo} \le x \le \val{xSlowHi}$,
$\val{ySlowLo} \le y \le \val{ySlowHi}$. The three fast shots are counted in the
lower spot $\val{xFastLo} \le x \le \val{xFastHi}$,
$\val{yFastLo} \le y \le \val{yFastHi}$. The player may place at most
$\val{maxCharges}$ charges, of either sign, inside
$\val{xRegLo} \le x \le \val{xRegHi}$, $\val{yRegLo} \le y \le \val{yRegHi}$,
and no magnets. A flight ends when the particle touches a charge, leaves the
world, enters its detector or runs out of time
(\xsecref{foundations}{sec:board}); $t_{\max} = \val{tmax}$.

The reference solution is six charges, all of the particle's sign:
$Q_1 = \val{Q1}$ at $(\val{x1}, \val{y1})$,
$Q_2 = \val{Q2}$ at $(\val{x2}, \val{y2})$,
$Q_3 = \val{Q3}$ at $(\val{x3}, \val{y3})$,
$Q_4 = \val{Q4}$ at $(\val{x4}, \val{y4})$,
$Q_5 = \val{Q5}$ at $(\val{x5}, \val{y5})$ and
$Q_6 = \val{Q6}$ at $(\val{x6}, \val{y6})$.
The first pair brings every ray through the gate. The four behind it send the
slow rays up into the upper spot and the fast rays down into the lower spot.
This document integrates the six rays, and the same rays with only the first
pair, by the fixed-step method of \levelref{around_the_wall}.

\section{Setup}\label{sec:setup}

The level file, field by field (the appendix reproduces it):
\setuptable

In symbols: $c = \val{c}$, the particle has radius $a = \val{aRad}$, and every
charge has radius $R = \val{R}$. The slow launch has
\[
  \gamma_0 = \val{gammaSlow}, \qquad \abs{\vect{p}_0} = \val{pSlow}, \qquad
  v_0 = \val{vSlow},
\]
and the fast launch has $\gamma_0 = \val{gammaFast}$,
$\abs{\vect{p}_0} = \val{pFast}$ and $v_0 = \val{vFast}$. The world is
$-1 \le x \le \val{xWorld}$, $-1 \le y \le \val{yWorld}$. The couplings are
$\kappa_1 = \val{k1}$, $\kappa_2 = \val{k2}$, $\kappa_3 = \val{k3}$,
$\kappa_4 = \val{k4}$, $\kappa_5 = \val{k5}$ and $\kappa_6 = \val{k6}$, in the
file's order, so the reference flight has $N = \val{nCharges}$ centres. The
lens placement keeps the first two and has $N = \val{nLens}$.

\begin{figure}[htbp]
\centering
\begin{tikzpicture}
\begin{axis}[axis equal image, width=0.92\linewidth, xmin=-1.5, xmax=42.5, ymin=-1.5, ymax=21.5,
  xtick={0,5,...,40}, ytick={0,5,...,20}, tick label style={font=\small},
  xlabel={$x$}, ylabel={$y$}, label style={font=\small}]
  \draw[black!45, dashed] (-1,-1) rectangle (41,21);
  \draw[black!50] (0,0) rectangle (40,20);
  \draw[black!35, dashed] (\rawval{xRegLo},\rawval{yRegLo}) rectangle (\rawval{xRegHi},\rawval{yRegHi});
  \fill[orange!70, fill opacity=0.25] (\rawval{xGlo},\rawval{yGlo}) rectangle (\rawval{xGhi},\rawval{yGhi});
  \draw[orange!70!black] (\rawval{xGlo},\rawval{yGlo}) rectangle (\rawval{xGhi},\rawval{yGhi});
  \fill[green!50!black, fill opacity=0.18] (\rawval{xSlowLo},\rawval{ySlowLo}) rectangle (\rawval{xSlowHi},\rawval{ySlowHi});
  \draw[green!40!black] (\rawval{xSlowLo},\rawval{ySlowLo}) rectangle (\rawval{xSlowHi},\rawval{ySlowHi});
  \fill[green!50!black, fill opacity=0.18] (\rawval{xFastLo},\rawval{yFastLo}) rectangle (\rawval{xFastHi},\rawval{yFastHi});
  \draw[green!40!black] (\rawval{xFastLo},\rawval{yFastLo}) rectangle (\rawval{xFastHi},\rawval{yFastHi});
  \fill[red!70!black, fill opacity=0.4] (\rawval{x1},\rawval{y1}) circle[radius=\rawval{R}];
  \fill[red!70!black, fill opacity=0.4] (\rawval{x2},\rawval{y2}) circle[radius=\rawval{R}];
  \fill[red!70!black, fill opacity=0.4] (\rawval{x3},\rawval{y3}) circle[radius=\rawval{R}];
  \fill[red!70!black, fill opacity=0.4] (\rawval{x4},\rawval{y4}) circle[radius=\rawval{R}];
  \fill[red!70!black, fill opacity=0.4] (\rawval{x5},\rawval{y5}) circle[radius=\rawval{R}];
  \fill[red!70!black, fill opacity=0.4] (\rawval{x6},\rawval{y6}) circle[radius=\rawval{R}];
  \addplot[orange!80!black] table {pathfastdown.dat};
  \addplot[purple!70!black] table {pathfastup.dat};
  \addplot[very thick, red!75!black] table {pathfastaxis.dat};
  \addplot[blue!70] table {pathslowdown.dat};
  \addplot[cyan!70!blue] table {pathslowup.dat};
  \addplot[very thick, blue!55!black] table {pathslowaxis.dat};
  \draw[-{Stealth[length=5pt]}, thick] (\rawval{xLaunch},\rawval{yLaunch}) -- (\rawval{arrowX},\rawval{arrowY});
  \fill (\rawval{xLaunch},\rawval{yLaunch}) circle[radius=0.08];
  \node[font=\scriptsize, orange!60!black, anchor=south] at (19,11.15) {gate};
  \node[font=\scriptsize, green!30!black, anchor=west] at (40.2,14) {slow};
  \node[font=\scriptsize, green!30!black, anchor=west] at (40.2,6) {fast};
\end{axis}
\end{tikzpicture}
\caption{The six reference rays. Blue is $T = \val{Tslow}$: the three rays meet
the left face of the upper spot, and the cyan ray is the one launched upward.
Red and orange are $T = \val{Tfast}$: the thick red ray meets the top of the
lower spot, and the orange and purple rays meet its left face. Every ray has
passed the gate. The discs are the six charges.}
\label{fig:board}
\end{figure}

\section{Principles and model}\label{sec:model}

\begin{principles}
  \item The field of a charged sphere and the force on a non-polarizable particle,
    \xeqref{foundations}{eq:coulomb} and \xeqref{foundations}{eq:sphereforce}.
  \item The relativistic equation of motion \xeqref{foundations}{eq:eom}, the launch
    \xeqref{foundations}{eq:launch} and the energy \xeqref{foundations}{eq:energy}. The
    motion stays in the plane (\xsecref{foundations}{sec:plane}).
  \item The board's rules (\xsecref{foundations}{sec:board}), and the gate
    \xeqref{two_stages}{eq:gate}. The detector counts only after the gate has been
    entered. This level adds no direction cone and no kinetic window.
  \item The classical Runge--Kutta method \xeqref{foundations}{eq:rk4} and the fixed-step
    flight of \xsecref{around_the_wall}{sec:method}. With two or more centres the orbit
    equation does not separate (\xeqref{slingshot}{eq:noseparation}).
  \item The transverse kick of a straight flyby scales as $1/T$
    (\xeqref{thomson_crt}{eq:impulse}). That is why the slower rays can be steered
    onto a different spot. The paths here are not that flyby, so the landing is the
    integral of the force.
  \item The radiated power \xeqref{foundations}{eq:larmor}, with the acceleration
    \xeqref{first_bend}{eq:acceleration}, summed by Simpson's rule
    \xeqref{foundations}{eq:simpson}.
\end{principles}

The level has no magnets, so $\vect{B} = \vect{0}$. The force is
\xeqref{around_the_wall}{eq:force} with $N = \val{nCharges}$ on the reference,
the charges in the file's order. The state $(x, y, p_x, p_y)$ obeys
\xeqref{foundations}{eq:eom} with that force until the first ending event. The
potential in \xeqref{around_the_wall}{eq:energy} is the same superposition. A
gate entry is not an ending: the flight continues, and only a later detector
entry is kept, and only when the gate time is earlier
(\xeqref{two_stages}{eq:gate}, \physref{6.2}).

\begin{assumption}{The charges are held in place}\label{as:fixed}
None of the charges recoils. That is a rule of the game
(\xsecref{foundations}{sec:radiation}).
\end{assumption}

\begin{assumption}{No radiation}\label{as:radiation}
The radiated energy is not taken from the particle. Over the twelve flights of
this document the largest fraction of the launch energy is $\val{radFracMax}$
(§\ref{sec:hand}). The game makes the same approximation and reports the same
integral (\physref{8}).
\end{assumption}

\begin{assumption}{A point particle}\label{as:point}
$a = \val{aRad}$, and the particle does not polarize, so
\xeqref{foundations}{eq:sphereforce} is exact. The closest approach in the
twelve flights is $\val{rCloseMin}$, on the fast downward reference ray, a
clearance of $\val{clearMin}$ outside the sphere. The samples that locate it are
at most $\val{sampleArc}$ apart (§\ref{sec:method}), well below that clearance.
\end{assumption}

\begin{assumption}{The step resolves the gate and the spots}\label{as:step}
The reported step is $h = \val{H}$. Halving it estimates the largest error of an
end state as $\val{errMax}$ (§\ref{sec:hand}). That is far below the height of
either spot and far below the height of the gate. The tolerance used in
§\ref{sec:comparison} is the larger of $10^{-9}$ and ten times this estimate.
\end{assumption}

\section{Derivation}\label{sec:derivation}

\subsection{The gate}\label{sec:gate}

The rule is \xeqref{two_stages}{eq:gate}. This level lists one gate and gives it
no condition on the velocity or the energy. If $t_G$ is the first time the
position lies in the gate and $t_D$ the first time it lies in that shot's
detector, the detector counts only when $t_G$ exists and $t_G < t_D$. An entry
of the detector with the gate still ahead is a skipped gate, the first in the
list. Leaving the world, or touching a charge, ends the flight on its own
outcome whether or not the gate was entered. On a path that has already curved,
$t_G$ is not a root of a separated orbit: it is the crossing found by
§\ref{sec:method}, and the flight is then continued from that step.

\subsection{Why the slower ray can be sorted}\label{sec:sort}

\xeqref{thomson_crt}{eq:impulse} is the angle of a straight flyby past one
charge, from $-\infty$ to $\infty$: it is proportional to $1/T$. The same
charges therefore turn a slower ray through a larger angle. The reference is
not that flyby. Two charges stand in front of the gate and four behind it, and
the rays leave the straight line at once. The formula is used here only for the
direction of the sort. The landings are the integral.

On the two axis rays the launch is along $+x$, so the angle at the gate is the
whole turn. The slow ray has turned by $\val{turnSlow}^\circ$ and the fast ray
by $\val{turnFast}^\circ$. The turns are in the ratio $\val{turnRatio}$; the
energies are in the ratio $\val{ratioT}$. The products of turn and energy are
$\val{kickSlow}$ and $\val{kickFast}$. A single straight flyby would make those
products equal. They are not. The slower ray has nevertheless turned further,
and §\ref{sec:hand} finds it in the upper spot while the faster ray meets the
lower one.

\subsection{Why the flight is integrated}\label{sec:noseparation}

With one charge, angular momentum closes the motion into the orbit of
\levelref{first_bend}. Two fixed centres already fail to separate in the
relativistic problem: the term \xeqref{slingshot}{eq:noseparation} does not
split. Six centres give that obstruction nothing that would remove it. Each ray
is the initial-value problem for \xeqref{around_the_wall}{eq:force}, and
§\ref{sec:method} computes it to a stated accuracy.

\subsection{The method}\label{sec:method}

The state is advanced by the method of \xsecref{around_the_wall}{sec:method}:
the classical Runge--Kutta method \xeqref{foundations}{eq:rk4} at the fixed step
$h = \val{H}$, a crossing on the cubic Hermite interpolant
\xeqref{slingshot}{eq:hermite}, $\val{nSample}$ samples and $\val{nBisect}$
bisections, and the enlarged-box test. A gate entry is recorded and the step
continues. A collision, the world's edge or the detector ends the flight, in
that order (\xsecref{foundations}{sec:board}). The face coordinate is then set
to the face value. If the ending is the detector and the gate has not been
entered, the outcome is a skipped gate. The error of the reported step is a
fifteenth of its difference from the step $2h$
(\xsecref{foundations}{sec:numerics}).

On each step the interpolant is also sampled at nine fractions, up to the
ending when the step contains one. Those samples give the closest approach and
the range of $y$. A point between them lies within $\val{sampleArc}$ of a
sample. The end state is held to the step-halving estimate.

\section{The reference solution by hand}\label{sec:hand}

\begin{procedure}
  \item\label{st:launch} \emph{The launch} (\xeqref{foundations}{eq:launch}).
    The slow rays have $\gamma_0 = \val{gammaSlow}$ and
    $\abs{\vect{p}_0} = \val{pSlow}$, the fast rays $\gamma_0 = \val{gammaFast}$
    and $\abs{\vect{p}_0} = \val{pFast}$. The aims are $\val{angSlowDown}^\circ$,
    $\val{angSlowAxis}^\circ$ and $\val{angSlowUp}^\circ$.
  \item\label{st:energy} \emph{The energy.} At the common launch point,
    $V = \val{RefSlowAxisVzero}$. Then $W - mc^2$ is $\val{RefSlowAxisWbind}$ on
    a slow ray and $\val{RefFastAxisWbind}$ on a fast ray. Both are positive, so
    the world's edge is open and the integration decides who reaches a spot.
  \item\label{st:gate} \emph{The gate} \xeqref{two_stages}{eq:gate}, on the slow
    axis ray. The crossing is the \val{RefSlowAxisGface} face at
    $(\val{RefSlowAxisGx}, \val{RefSlowAxisGy})$, at
    $t = \val{RefSlowAxisGt}$, with the velocity at $\val{RefSlowAxisGang}^\circ$
    to $+x$. The residual is $\val{RefSlowAxisGres}$.

    The other five reference rays also enter the gate, all before any detector.
    The slow downward ray enters the \val{RefSlowDownGface} face at
    $(\val{RefSlowDownGx}, \val{RefSlowDownGy})$, and the slow upward ray the
    \val{RefSlowUpGface} face at $(\val{RefSlowUpGx}, \val{RefSlowUpGy})$.

    The fast downward ray enters the \val{RefFastDownGface} face at
    $(\val{RefFastDownGx}, \val{RefFastDownGy})$, the fast axis ray the
    \val{RefFastAxisGface} face at $(\val{RefFastAxisGx}, \val{RefFastAxisGy})$,
    and the fast upward ray the \val{RefFastUpGface} face at
    $(\val{RefFastUpGx}, \val{RefFastUpGy})$.
  \item\label{st:flight} \emph{The slow axis flight} after the gate, at
    $h = \val{H}$. The states every five time units, and the drift of the energy:
    \begin{center}\small\val{states}\end{center}
    The ray, which crossed the gate heading down, is turned upward. It enters
    the \val{RefSlowAxisFace} face of the upper spot at
    $(\val{RefSlowAxisX}, \val{RefSlowAxisY})$, at $t = \val{RefSlowAxisT}$,
    after $\val{RefSlowAxisSteps}$ steps. The momentum is
    $(\val{RefSlowAxisPx}, \val{RefSlowAxisPy})$, at
    $\val{RefSlowAxisAng}^\circ$ to $+x$. The ordinate lies between
    $\val{ySlowLo}$ and $\val{ySlowHi}$.
  \item\label{st:others} \emph{The other five arrivals.} Each has passed the
    gate, and each is Arrived.

    The slow downward ray meets the \val{RefSlowDownFace} face at
    $(\val{RefSlowDownX}, \val{RefSlowDownY})$, at $t = \val{RefSlowDownT}$,
    heading $\val{RefSlowDownAng}^\circ$.

    The slow upward ray climbs to $\val{RefSlowUpYmax}$ and meets the
    \val{RefSlowUpFace} face at $(\val{RefSlowUpX}, \val{RefSlowUpY})$, at
    $t = \val{RefSlowUpT}$, heading $\val{RefSlowUpAng}^\circ$.

    The fast downward ray meets the \val{RefFastDownFace} face of the lower
    spot at $(\val{RefFastDownX}, \val{RefFastDownY})$, at
    $t = \val{RefFastDownT}$, heading $\val{RefFastDownAng}^\circ$.

    The fast axis ray meets the \val{RefFastAxisFace} face at
    $(\val{RefFastAxisX}, \val{RefFastAxisY})$, at $t = \val{RefFastAxisT}$,
    heading $\val{RefFastAxisAng}^\circ$. The ordinate is the top of the spot.

    The fast upward ray meets the \val{RefFastUpFace} face at
    $(\val{RefFastUpX}, \val{RefFastUpY})$, at $t = \val{RefFastUpT}$, heading
    $\val{RefFastUpAng}^\circ$.
  \item\label{st:check} \emph{The checks.} Along the slow axis flight the energy
    stays within $\val{RefSlowAxisWdev}$ of $W$. Halving the step estimates the
    six reference end states as $\val{RefSlowDownErr}$, $\val{RefSlowAxisErr}$,
    $\val{RefSlowUpErr}$, $\val{RefFastDownErr}$, $\val{RefFastAxisErr}$ and
    $\val{RefFastUpErr}$. The same outcomes appear at $h$ and at $2h$. The
    forward momentum stays positive on every reference ray; the smallest value
    is $\val{RefSlowUpPxMin}$, on the slow upward ray. The closest approach is
    $\val{RefFastDownRclose}$.
  \item\label{st:rad} \emph{The radiation}, on the two axis rays: Simpson's rule
    on the steps of $h/2$. The slow axis ray radiates $\val{RefSlowAxisRad}$ and
    the fast axis ray $\val{RefFastAxisRad}$. The relative differences from the
    same sums at step $h$, divided by fifteen, are $\val{RefSlowAxisRadErr}$ and
    $\val{RefFastAxisRadErr}$.
\end{procedure}

\begin{result}[Six rays, two spots]
Every reference ray enters the gate and then its own spot. The three slow rays
meet the left face of the upper spot, at ordinates $\val{RefSlowDownY}$,
$\val{RefSlowAxisY}$ and $\val{RefSlowUpY}$. The fast downward and fast upward
rays meet the left face of the lower spot, at $\val{RefFastDownY}$ and
$\val{RefFastUpY}$. The fast axis ray meets the top of that spot, at
$x = \val{RefFastAxisX}$. The level is solved.
\end{result}

By hand, the slow axis ray is $\val{RefSlowAxisSteps}$ steps of four evaluations
of the force, and the other five rays are the same method. Across the twelve
flights the largest step-halving estimate is $\val{errMax}$.

\section{Comparison with the game}\label{sec:comparison}

\comparisontable

Thirty-eight comparisons are in the table. For each reference ray the outcome,
the time and the end position are compared. For the first two charges alone,
marked L, the outcome and the time are compared; the end positions are in
§\ref{sec:others}. The radiated energy is compared on the two reference axis
rays. The tolerance on a state is $\val{tolFlight}$. Ten times the largest
step-halving estimate $\val{errMax}$ lies above $10^{-9}$, so the tolerance is
that product rather than the floor. The radiation tolerance is $10^{-4}$,
relative. The state tolerance and the radiation tolerance were fixed from these
estimates before the comparison.

The largest absolute difference is the ordinate of the fast downward ray,
$\val{cmpRefFastDownYDiff}$. The next is the ordinate of the slow upward ray,
$\val{cmpRefSlowUpYDiff}$. Both are smaller than $\val{errMax}$ and a fraction
of the tolerance. The game's verified end states are good to about $10^{-10}$
(\xsecref{foundations}{sec:verification}).

The abscissae of the five rays that meet a detector's left face agree exactly.
The fast axis ray meets the top face, so its abscissa is not that coordinate;
it differs by $\val{cmpRefFastAxisXDiff}$. The quadrature uncertainties are
$\val{RefSlowAxisRadErr}$ and $\val{RefFastAxisRadErr}$. The larger relative
difference of the two is the fast axis ray, $\val{cmpRefFastAxisRadDiff}$.

\section{Other solutions}\label{sec:others}

The first two charges, and nothing behind the gate, do not solve the level.
This is the lens placement in \code{placements.json}: $Q_1$ and $Q_2$ only.
The method of §\ref{sec:method} integrates all six rays at $h = \val{H}$ and at
$2h$. Every one of them still enters the gate. What happens afterwards is not
one outcome.
\begin{enumerate}[label=\arabic*.]
  \item The three slow rays leave through the right face of the world, under
    the lower spot. The downward ray leaves at
    $(\val{LensSlowDownX}, \val{LensSlowDownY})$, at $t = \val{LensSlowDownT}$.
    The axis ray leaves at $(\val{LensSlowAxisX}, \val{LensSlowAxisY})$, at
    $t = \val{LensSlowAxisT}$. The upward ray leaves at
    $(\val{LensSlowUpX}, \val{LensSlowUpY})$, at $t = \val{LensSlowUpT}$.
    None of these ordinates reaches $\val{ySlowLo}$.
  \item The fast downward ray does arrive. It meets the
    \val{LensFastDownFace} face of the lower spot at
    $(\val{LensFastDownX}, \val{LensFastDownY})$, at $t = \val{LensFastDownT}$.
  \item The fast axis ray also arrives. It meets the \val{LensFastAxisFace}
    face at $(\val{LensFastAxisX}, \val{LensFastAxisY})$, at
    $t = \val{LensFastAxisT}$.
  \item The fast upward ray misses the lower spot. It leaves through the right
    face of the world at $(\val{LensFastUpX}, \val{LensFastUpY})$, at
    $t = \val{LensFastUpT}$. The ordinate is above $\val{yFastHi}$.
\end{enumerate}

The pair in front of the gate is enough to pass all six rays through it, and
enough to land the fast downward ray and the fast axis ray. It does not lift
the slow rays, and the fast upward ray stays just above the lower spot. The
four charges behind the gate are what turn the slow rays up into the upper spot
and bring that last fast ray down into the lower one. The two fast rays that
already arrived still arrive, on the faces reported in §\ref{sec:hand}.
Section~\ref{sec:comparison} records the outcomes.

\section{What the game computes}\label{sec:engine}

The game integrates \xeqref{foundations}{eq:eom} with its eighth-order
Runge--Kutta method at the tolerances $10^{-10}$ and $10^{-12}$ (\physref{5}).
It finds each crossing on its interpolant (\physref{6}), continues through the
gate, and counts the detector only after that entry (\physref{6.2}). It
verifies the outcome by comparing its two runs (\physref{7}). The model is this
document's, including \ref{as:radiation} (\physref{8}). The agreement in
§\ref{sec:comparison} is between the fixed-step flight of §\ref{sec:method},
the gate of §\ref{sec:gate}, and that integration.

\section{Further reading}\label{sec:reading}

\begin{itemize}
  \item \levelref{two_stages}: the gate \xeqref{two_stages}{eq:gate}, applied
    here to six rays.
  \item \levelref{thomson_crt}: the kick proportional to $1/T$
    (\xeqref{thomson_crt}{eq:impulse}), which says why the slower rays can be
    placed on the other spot.
  \item \levelref{around_the_wall}: the force of several centres and the
    fixed-step method, used here at $N = \val{nCharges}$.
  \item \levelref{slingshot}: why two fixed centres have no closed relativistic
    orbit, and the Hermite crossing.
  \item \levelref{first_bend}: the acceleration that enters the radiated power.
  \item \xsecref{foundations}{sec:board} and \xsecref{foundations}{sec:numerics}:
    the order of events, and the step-halving estimate.
\end{itemize}

\levelfileappendix

\end{document}
